\scanpage{50}
\mspage{51}

\deleted{L'exemple le plus important de produit tensoriel hilbertien
est \unsure{donné|fourni} par le

\uline{\unsure{Th|Prop}} Soit $E$ un espace de Hilbert,
\deleted{soit $\mu$ une mesure
sur un \illegible{2}}
positive sur un espace loc\textsuperscript{t} compact $M$, considérons,
l'espace de Hilbert $L^2(\mu)$, et l'espace de Hilbert}

\uline{Exemple} Soit $M$ un espace loc\textsuperscript{t} cpt muni d'une
mesure positive $\mu$, considérons l'espace de Hilbert
$L^2(\mu)$ (où $(\varphi,\psi)=\int\varphi\overline\psi\,d\mu$).
Pour tout espace de Banach $E$, $L^2_E(\mu)$ désigne l'espace des
fonctions sur $M$, à valeurs dans $E$, de \unsure{puissance deuxième}
intégrables, muni de la norme
\[
(9)\qquad
\|f\|_2=\left(\int \deleted{\unsure{f(t)}}\|f(t)\|_E^2\,d\mu(t)\right)^{1/2}
\]
qui en fait un espace de Banach $(\qquad)$. Si
$E$ est un espace de Hilbert, la norme $(9)$ peut
aussi être définie par un produit scalaire
\[
(10)\qquad (f,g)=\int\langle f(t),\overline g(t)\rangle\,d\mu(t)
\]
donc $L^2_E(\mu)$ est alors un espace de Hilbert.

On a une application bilinéaire naturelle
$(\varphi,a)\longrightarrow\varphi.a$ de $L^2\deleted{\unsure{(\mu)}}\times E$ dans $L^2_E$,
d'où une application linéaire \inserted{$u\to\widetilde u$}
$L^2\otimes E\longrightarrow L^2_E$, dont
l'image est d'ailleurs dense dans $L^2_E$. \deleted{\unsure{Il en}}
\deleted{\unsure{qui|par}} C'est un isomorphisme pour les structures
préhilbertiennes, car \inserted{on a évidemment}
$(\varphi.a,\psi.b)=(\varphi,\psi)(a,b)$,
i.e
\[
(\widetilde{\varphi\otimes a},\widetilde{\psi\otimes b})
=(\varphi\otimes a,\psi\otimes b),
\]
d'où \unsure{linéairement|linéarité} $(\widetilde u,\widetilde v)=(u,v)$
pour tous $u,v\in L^2\otimes E$.
